% !PRINTED 44
\weberpage{44}
%<<CONT
\gesperrt{Abelschen Theorems} und soll auch in dieser Weise hier aufgefa\ss t und abgeleitet werden\footnote{Dieses weitumfassende Theorem findet sich in gro\ss artiger Einfachheit abgeleitet in einem kaum zwei Seiten umfassenden Aufsatze im vierten Bande von \gesperrt{Crelles} Journal \glqq D\'emonstration d'une propri\'et\'e g\'en\'erale d'une certaine classe de fonctions transcendentes\grqq{}; Oeuvres compl\`etes de N. H. \gesperrt{Abel}. Nouvelle \'edition, T.\,I, p.\,515. Ausf\"uhrliche Darstellung und Anwendungen in der nachgelassenen gro\ss en Abhandlung: \glqq M\'emoire sur une propri\'et\'e g\'en\'erale d'une classe tr\`es \'etendue de fonctions transcendentes\grqq{}.}. Dem Beweise schicken wir folgenden elementaren algebraischen Satz voraus:

\gesperrt{Ist $F(x)$ eine ganze rationale Funktion $n$ten Grades ohne mehrfache Faktoren, $F'(x)$ ihre Derivierte, sind ferner $x_1, x_2, \ldots, x_n$ die Wurzeln der Gleichung $F(x) = 0$ und $\varphi(x)$ eine ganze Funktion von $x$, deren Grad h\"ochstens $= n-2$, so ist}
\weq{2}{\frac{\varphi(x_1)}{F'(x_1)} + \frac{\varphi(x_2)}{F'(x_2)} + \cdots \frac{\varphi(x_n)}{F'(x_n)} = 0.} % [sic] no plus sign is printed after the ellipsis

Der Beweis dieses Satzes ergibt sich unmittelbar aus der Zerlegung des rationalen Bruches
\[
\frac{x\,\varphi(x)}{F(x)}
\]
in Partialbr\"uche
\[
\frac{x_1\varphi(x_1)}{F'(x_1)(x-x_1)} + \frac{x_2\varphi(x_2)}{F'(x_2)(x-x_2)} + \cdots + \frac{x_n\varphi(x_n)}{F'(x_n)(x-x_n)},
\]
wenn darin $x = 0$ gesetzt wird [Bd.\,I, \S\ 15 (9)].

Es sollen nun $P$, $Q$ ganze Funktionen von $x$ von den Graden $m$ und $m-2$ bedeuten, und wir fragen nach den Werten von $x$, $\sqrt{f(x)}$, f\"ur die die Funktion
\weq{3}{P + Q\sqrt{f(x)}}
verschwindet. Solcher Wertpaare (Punkte) gibt es $2m$, und zwar findet man die Werte von $x$ als Wurzeln der Gleichung $2m$ten Grades:
\weq{4}{F(x) = P^2 - Q^2 f(x) = 0,}
und die zugeh\"origen Werte von $\sqrt{f(x)}$ aus
\weq{5}{P + Q\sqrt{f(x)} = 0.}

Wir nehmen nun an, die Koeffizienten in den Funktionen $P$, $Q$ seien ver\"anderlich, entweder unabh\"angige Ver\"anderliche oder in irgend einer Weise von anderen Ver\"anderlichen abh\"angig.
